\documentclass[12pt]{article}
\usepackage[T1]{fontenc}
\usepackage[utf8]{inputenc}
\usepackage{lmodern}
\usepackage{textcomp}
\usepackage{enumitem}
\usepackage{geometry}
\geometry{margin=1in}
\begin{document}
\begin{center}
Answer Key - Sociology - Graduate Level Assessment 17 \
\small (Answers are ordered exactly as in the question paper.)
\end{center}

\section*{Multiple Choice}
\begin{enumerate}[label=\arabic*.]
\item \textbf{Answer:} D
\\ \textit{Options:} A) the privatization of public services by the Conservative government; B) the lifestyle practice of shopping in peer groups; C) the form of tuberculosis suffered by those who collect stamps; D) the provision of health, housing, and education services by the state
\\ \textit{Note:} The correct answer is based on the concept of 'collective consumption' as described by Manuel Castells, which emphasizes the role of the state in providing essential services such as health, housing, and education to ensure social welfare and address collective needs in society.
\item \textbf{Answer:} D
\\ \textit{Options:} A) They are cultures as well as churches.; B) They practice separation of church and state.; C) They exclude women as clergy.; D) They are monotheistic.
\\ \textit{Note:} The correct answer is that these religious organizations are monotheistic because they all believe in the existence of a single, all-powerful deity, which is a defining characteristic of monotheistic faiths.
\item \textbf{Answer:} D
\\ \textit{Options:} A) the ruling class, or bourgeoisie, who exploit the proletariat; B) a capitalist class dependent on property ownership and advantaged life chances; C) an alignment of classes with shared interests but no state power; D) a state elite whose members are drawn overwhelmingly from a power bloc
\\ \textit{Note:} The term 'power elite' introduced by Scott (1991) accurately describes a state elite composed primarily of individuals from a power bloc, highlighting the concentration of power among a select group that influences political, economic, and social decisions.
\item \textbf{Answer:} D
\\ \textit{Options:} A) the white population did not believe in liberty, equality, and democracy; B) they wanted to retain a strong sense of their original ethnic identity; C) they were not prepared to leave the Southern states and move to North America; D) the promise of citizenship was contradicted by continued discrimination
\\ \textit{Note:} Warner and Myrdal argued that the persistent discrimination faced by black former slaves undermined the promise of citizenship, preventing their full integration into the 'ethnic melting pot' of North American cities.
\item \textbf{Answer:} A
\\ \textit{Options:} A) a religious organization that claims total spiritual authority over its members; B) a church organized around voluntary rather than compulsory membership; C) a sect or cult with a very small following; D) a hierarchy of priests or other spiritual leaders
\\ \textit{Note:} An ecclesia is characterized by its institutional role as a dominant religious organization that integrates fully with the state and claims comprehensive spiritual authority over its adherents, thereby distinguishing it from other forms of religious groups.
\end{enumerate}

\section*{True / False}
\begin{enumerate}[label=\arabic*.]
\setcounter{enumi}{5}
\item \textbf{Answer:} True
\\ \textit{Options:} A) True; B) False
\\ \textit{Note:} The office is primarily a workplace setting focused on organizational function and economic activity, whereas the household, nation state, and global village are recognized as distinct levels of social organization that encompass broader social relationships and structures.
\item \textbf{Answer:} False
\\ \textit{Options:} A) True; B) False
\\ \textit{Note:} Mead's 'generalized other' refers to the internalized attitudes, expectations, and viewpoints of the broader society that individuals use to guide their behavior, rather than being defined as a specific significant figure in childhood.
\item \textbf{Answer:} False
\\ \textit{Options:} A) True; B) False
\\ \textit{Note:} Booth's study indicated that approximately 30.7\% of Londoners were living in poverty, not 27.50\%, making the statement false.
\item \textbf{Answer:} True
\\ \textit{Options:} A) True; B) False
\\ \textit{Note:} Marx's concept of class polarization highlights the structural economic disparities created by capitalism, where the concentration of wealth among the bourgeoisie exacerbates the divide between social classes, ultimately fostering a heightened awareness and solidarity among the proletariat regarding their shared struggles.
\item \textbf{Answer:} True
\\ \textit{Options:} A) True; B) False
\\ \textit{Note:} The statement is true because the rise of industrialization in the nineteenth century fueled military advancements, while the competition among rival empires for global dominance intensified conflicts, making war between nation-states increasingly likely.
\end{enumerate}

\section*{Long Form Response}
\begin{enumerate}[label=\arabic*.]
\setcounter{enumi}{10}
\item \textbf{Answer:} Informal negative sanctions play a significant role in shaping social behavior by influencing individuals to conform to societal norms. In the case of children who are teased for thumb-sucking upon entering kindergarten, the ridicule they experience serves as an informal negative sanction that discourages this behavior. Such teasing can lead to feelings of shame and social exclusion, compelling children to abandon thumb-sucking in order to gain acceptance from their peers. This process highlights the power of social interactions in regulating behavior, as children learn to align their actions with the expectations of their social environment to avoid negative feedback. Ultimately, informal negative sanctions like teasing not only reinforce behavioral norms but also contribute to the internalization of societal values among young individuals.
\\ \textit{Note:} correct because it illustrates how informal negative sanctions, such as teasing, effectively discourage behaviors that deviate from social norms by instilling feelings of shame and the desire for social acceptance among children.
\item \textbf{Answer:} Nineteenth-century scientific theories often sought naturalistic explanations of race by emphasizing biological determinism and the idea that physical characteristics, such as height, were indicative of inherent racial qualities. Proponents like Samuel Morton argued that cranial capacity correlated with intelligence, while others posited that racial differences in stature reflected evolutionary superiority. However, these discussions frequently overlooked the socio-cultural factors influencing human development and the interaction between environment and genetics. Additionally, the focus on physical traits like height tended to reinforce social hierarchies rather than acknowledge the complexities of race as a social construct. This reductionist approach ultimately oversimplified human diversity and ignored the significant role of historical and contextual influences.
\\ \textit{Note:} correct because it identifies the emphasis on biological determinism and physical characteristics in nineteenth-century scientific theories of race while also highlighting the critical oversight of socio-cultural influences in these discussions.
\end{enumerate}

\vfill
\noindent\textit{End of Answer Key}
\end{document}